% PayBench methods paper — arXiv preprint source (derived from paybench/methodology/methodology.md v1.1)
% Build: pdflatex payhelm-methods.tex  (twice, for refs)
% Perimeter discipline (methodology §7, F14): this PUBLIC artefact carries NO real-rail
% ranking and NO per-rail finality numbers. Calibration parameters live in the repository.
% Watermark (cross-LLM review): the title states "Calibrated Mock Harness" so the work cannot
% be screenshotted as a real-rail leaderboard.
\documentclass[11pt]{article}
\usepackage[margin=1in]{geometry}
\usepackage{amsmath}
\usepackage{booktabs}
\usepackage{hyperref}
\hypersetup{colorlinks=true, linkcolor=blue, urlcolor=blue, citecolor=blue}

\title{PayBench: A Pre-Registered Methodology and Calibrated Mock Harness\\
for Agent-to-Agent Payment-Rail Settlement Finality}

% TODO(identity): author + affiliation are placeholders pending the Stage-1 vs Stage-2
% identity decision (research-framing vs founder-framing). Do NOT submit with these as-is.
\author{[AUTHOR — identity decision pending]\\ \texttt{[affiliation pending]}}
\date{Methodology v1.1 (frozen 2026-06-06). Preprint v1.}

\begin{document}
\maketitle

\begin{abstract}
\noindent\textbf{This preprint describes a pre-registered measurement methodology and a calibrated
\emph{mock} harness; no real-rail funds are moved and no production rail-by-rail ranking is derived
or published here.} PayBench is an evaluation methodology and benchmark for agent-to-agent (A2A)
payment rails. Its first dimension is \emph{settlement-finality}: the wall-clock time from payment
submission to the point at which an autonomous agent may rely on the payment for an irreversible
downstream action. We make three contributions. First, a \emph{reliance-level doctrine} for defining
finality across heterogeneous rails (deterministic BFT close, probabilistic rooting, optimistic-rollup
soft finality, hash-time-locked preimage release), pairing each rail's ecosystem-canonical reliance
level with an explicit \emph{trust/equivalence class} so that wall-clock comparisons never masquerade
as security-model equivalence. Second, a comparative statistical method---Bradley--Terry maximum-likelihood
estimation over pairwise finality races, complemented by $\mathrm{pass}@k = P(\text{finality} \le k)$
and Wilson lower-bound confidence intervals---with explicit caveats on its behaviour under cluster
separation. Third, a cryptographic \emph{pre-registration} protocol (OSF DOI $+$ a transparency-log
entry $+$ an independent timestamp anchor $+$ a signed source tag $+$ this preprint) that fixes the
design before any scored run, and that cleanly decouples \emph{methodology} pre-registration from the
later \emph{real-rail data} pre-registration. The methodology and a reproducible mock harness are
released as an open tool; publication of real-rail rankings is deferred.
\end{abstract}

\section{Introduction}
Agentic payment rails are proliferating---x402 on multiple chains, the Machine Payments Protocol (MPP)
on different settlement rails, mandate/authorization layers such as AP2, and others. An autonomous
agent transacting A2A must answer one question above all others before it acts on a payment:
\emph{when can I rely on this?} That is the settlement-finality question, and it gates every
downstream action the agent takes. PayBench exists to define and publish the canonical, comparative
metric for it.

The benchmark is deliberately \emph{comparative}, not absolute: it ranks rails against each other on a
dimension, with a stated confidence interval, under a protocol fixed in advance. The comparative
framing is what makes a Bradley--Terry method the right tool, and what keeps the claims defensible.

Two design commitments shape everything below. (i)~\emph{Pre-registration.} The rail set, per-rail
finality definitions, trial counts, metric definitions, the random seed, and the input fixture hashes
are committed---cryptographically---before any scored run, so that no element can later be accused of
having been chosen to flatter a rail. (ii)~\emph{A calibrated mock harness.} The released tool ships
with calibrated mock-data fixtures as its test harness, exercising the full measurement and statistics
pipeline against realistic distributions \emph{without} publishing a real-rail ranking. Running the
method against production rails and publishing rail-by-rail scoring is deferred.

\section{What PayBench measures}
PayBench answers three different \emph{kinds} of question and is explicit about which is which (a
\emph{hybrid-tiered} schema):
\begin{itemize}
  \item \textbf{Tier~1 --- benchmarked.} Dimensions PayBench runs trials on, pre-registers, and ranks
  with confidence intervals. Settlement-finality is the first; authorization latency and fee
  predictability are roadmap-committed.
  \item \textbf{Tier~2 --- factual lookup.} Live, queryable capability/reference data per rail (fees,
  supported assets, custody model, the auth/dispute/refund/sanctions capability bundle), stated as
  fact, not scored.
  \item \textbf{Tier~3 --- roadmap.} Dimensions named for transparency about direction, carrying no
  measurements.
\end{itemize}
The tiering is an honesty mechanism: measured, looked-up, and roadmapped claims stay visibly distinct,
so nothing reads as a result that is not one.

\section{Settlement-finality}\label{sec:finality}
\paragraph{Definition.} Settlement-finality is the wall-clock time from payment \emph{submission} to
the point at which the value transfer is irreversible on the rail---the moment the payee can rely on
the funds without risk of reversal.

\paragraph{The reliance-level doctrine.} Rails do not share a single cryptographic finality model, so
no single uniform threshold can be pinned across all of them. Instead PayBench pins each rail to its
\emph{ecosystem-canonical reliance level}---the point that rail's own protocol/documentation
designates as the level at which a builder may rely on the payment for an irreversible downstream
action. This is the correct doctrine for an \emph{agent-DX} benchmark: it measures the moment an
autonomous agent is told it can act.

\paragraph{The doctrine's sharp edge, named openly.} Because the canonical reliance level differs by
rail, the thresholds are not equi-conservative: an optimistic-rollup's canonical level is
\emph{optimistic} (soft finality), whereas a probabilistic L1's canonical level may be
\emph{conservative} (a rooted/finalized commitment rather than an optimistic one). A uniform-optimistic
or a uniform-irreversibility doctrine would each be internally symmetric but would each misrepresent at
least one rail's own builder guidance. PayBench takes the per-rail-canonical doctrine and \emph{names}
the asymmetry rather than hiding it. To make it impossible to miss, every rail carries a
\textbf{trust/equivalence-class} label (Table~\ref{tab:trust})---a centralised-sequencer soft-commit is
not the same security object as a decentralised-supermajority cryptographic root, even when both are
that rail's canonical reliance point. The comparison is wall-clock-to-reliance, not
security-model-to-security-model, and the method says so.\footnote{This doctrine was upheld $3$--$1$ by
an independent four-model adversarial review against the uniform-optimistic alternative, on the ground
that ranking a rail at its \emph{optimistic} commitment would publish a threshold that rail's own
guidance disavows for irreversible reliance---i.e.\ measurement invalidity.}

\begin{table}[t]
\centering
\caption{Trust/equivalence classes (illustrative of the doctrine; per-rail finality definitions and
calibration parameters are in the repository, not reproduced here---see \S\ref{sec:design}).}
\label{tab:trust}
\begin{tabular}{ll}
\toprule
Finality mechanism & Trust / equivalence class \\
\midrule
Optimistic-rollup soft finality (1 block) & Optimistic soft-commit (centralised sequencer promise) \\
Deterministic ledger/BFT close & Deterministic BFT \\
Probabilistic rooted commitment & Cryptographic root (decentralised supermajority) \\
Hash-time-locked preimage release & HTLC claim (possibly via hosted-operator ceremony) \\
\bottomrule
\end{tabular}
\end{table}

\paragraph{Unit.} Seconds, wall-clock, from request to the \emph{agent-perceived finality signal}---the
point at which the agent is told it may act, which per rail adapter may be an HTTP response, a webhook,
or a polled state transition (not necessarily the first acceptance response). This deliberately bundles
the rail's protocol/transport envelope with on-chain settlement, because that full round-trip is what
the agent actually experiences and waits on.

\paragraph{Coverage note.} Authorization/mandate layers that do not settle independently (their finality
reduces to an underlying rail) are excluded from the finality dimension and measured on an
authorization-latency dimension instead, where they are well-defined. Racing such a layer against its
own underlying rail would violate the independence the statistical model assumes.

\section{Statistical method}
\paragraph{Bradley--Terry MLE.} For each unordered rail pair we run repeated finality \emph{races}: in a
trial, both rails execute the same scenario against the same fixture and the rail reaching finality
first wins. Aggregating wins/losses across all trials and pairs, Bradley--Terry MLE produces a latent
strength per rail,
\[
  P(i \succ j) = \frac{e^{\beta_i}}{e^{\beta_i} + e^{\beta_j}},
\]
a single comparative ranking internally consistent across all pairwise comparisons.

\paragraph{Caveat---separation.} When rails fall into widely separated latency regimes, slower rails
essentially never beat faster ones and the MLE approaches complete separation; a symmetric smoothing
prior keeps it finite. In that regime the BT strength \emph{magnitudes} are an artefact of the prior
and the trial count---only the \emph{rank order} and any \emph{within-regime} differences are
interpretable. We therefore do not present fine-grained strengths as precise, and publish the raw
pairwise win/loss matrix so a reader can verify the ranking is data-driven, not prior-driven.
$\mathrm{pass}@k$ is the metric that actually discriminates in this regime.

\paragraph{$\mathrm{pass}@k$.} Operationally an agent needs an absolute answer: ``will this rail settle
within my SLA?'' We report $\mathrm{pass}@k = P(\text{finality} \le k\text{ seconds})$ for a frozen
grid $k \in \{2,3,5,10,15,20\}$~s. The grid is justified on \emph{external agent-SLA grounds}---$2$~s as
an interactive-HTTP responsiveness threshold, $3$/$5$~s as synchronous-call budgets, $10$~s as a
gateway/timeout boundary, $15$/$20$~s as batch/background-settlement budgets---not reverse-engineered
from the data. It spans the full observed finality range so it discriminates both within the
fast-settling and within the slow-settling regimes, with $10$~s as the separator. The grid is frozen
with the rest of the design and is not re-fit per run.

\paragraph{Honesty note on grid timing.} Calibration data was collected \emph{before} the methodology
freeze, so the grid was finalised with the calibrated distributions in hand; the \emph{mock} run
therefore validates the measurement/statistics \emph{pipeline} against known inputs rather than testing
an independent hypothesis. The grid stands on the external SLA rationale above, and the \emph{same
frozen grid carries unchanged to the eventual real-rail run}, where it is fixed before that data
exists. Pre-registration here prevents $p$-hacking on the real-rail run; it does not, and is not claimed
to, retroactively blind the designers to the calibration data.

\paragraph{Wilson lower bounds, and their scope.} All proportions (pairwise win rates and
$\mathrm{pass}@k$) are reported with Wilson score intervals, leading with the $95\%$ lower bound as the
conservative figure; the Wilson interval is well-behaved near $0$ and $1$, where a normal approximation
misleads. Two scope caveats: (i)~the intervals are \emph{pointwise} $95\%$, with no family-wise or
false-discovery correction across the set, as they are reported as independent descriptive metrics
rather than a joint hypothesis test; (ii)~on \emph{mock} fixtures a Wilson interval captures the
Monte-Carlo sampling error of the harness, not real-world network variance---which is a separately
disclosed limitation.

\section{Experimental design and the calibrated mock harness}\label{sec:design}
The settlement-finality benchmark races the rails that settle independently. For the present
configuration this is $C(5,2)=10$ unordered pairs at $500$ trials per pair $=5{,}000$ symmetric trials.
A single published master seed drives all stochastic fixture selection and trial ordering (stdlib
Mersenne Twister; no wall-clock, no OS entropy), and every fixture is content-addressed (sha256) so a
trial records, and the harness re-verifies, the exact fixture bytes it consumed. The harness is pure
standard library, removing cross-version RNG drift.

\paragraph{Calibrated mock fixtures.} Each fixture is a seeded log-normal population whose location and
scale parameters are derived from real reference data---rail documentation, public telemetry, and
first-party testnet/devnet/regtest observations. The harness is therefore \emph{meaningful}: it
exercises the full pipeline against realistic, first-party-anchored distributions.

\paragraph{Perimeter discipline---why no per-rail numbers appear here.} A seconds-valued,
fastest-to-slowest table of real first-party measurements \emph{is} a rail ranking, which the
publication posture defers. Accordingly this preprint reports \emph{no} per-rail finality medians and
\emph{no} ranking. The calibration parameters (log-space $\mu$, $\sigma$ per rail), full per-source
provenance, the harness, and the deterministic run report are released in the repository; this paper
documents the \emph{method}. The released mock run demonstrates the pipeline; it is not, and is not
presented as, a real-rail result.

\paragraph{Disclosed limitations (no silent caps).} The following are committed \emph{into} the frozen
method so a later real-rail run is measured against a prior that already names them:
(1)~quiet testnet/devnet/regtest conditions understate mainnet congestion tails, so the fixture scale
parameters are an optimistic tail estimate; (2)~one rail's fixture is a manual-path \emph{lower bound}
on its spec path; (3)~a testnet-calibrated rail may differ from mainnet in empirical latency despite a
shared consensus algorithm. Central tendencies are high-confidence first-party; tails are the disclosed
soft spot. For the real-rail run, parameter uncertainty in the scale estimate (a modest per-rail sample)
should additionally be bootstrapped and propagated, and deterministic-finality rails should use a
discrete (constant-plus-jitter) distribution rather than log-normal.

\section{Pre-registration protocol}
Pre-registration is the load-bearing artefact: it is what defends a comparative leaderboard against
selective-trial and threshold-shopping critique. We pre-register, before any scored run, the rail set,
per-rail finality definitions, trial counts, metric definitions (BT $+$ $\mathrm{pass}@k$ grid $+$
Wilson), the RNG seed, and the input fixture hashes. The concrete frozen values are committed as a
content manifest; the manifest's own hash is the single value the anchors below commit to.

\paragraph{Cryptographic grounding (defence-in-depth).} The manifest hash is committed via three
independent trust anchors plus two corroborating records: an OSF pre-registration resolving to a DOI;
a signature over the manifest logged to a public transparency log; an independent timestamp anchored to
a public blockchain; a signed source-control tag on the frozen commit (hardware-key signed); and this
preprint. The triad (DOI $+$ blockchain timestamp $+$ transparency-log entry) provides no single point
of trust; the transparency-log entry and timestamp proof are independently sufficient to establish the
freeze date even if the DOI host is unavailable. Bulk trial telemetry, when real-rail runs occur, is
signed with a Merkle root plus the same signing mechanism.

\paragraph{Methodology vs.\ data pre-registration.} This artefact pre-registers the \emph{measurement
method and the calibrated mock baseline}. The eventual real-rail run is a \emph{separate}
pre-registration: its specific production-fixture hashes and run configuration are committed, with the
same anchor stack, \emph{before} that run executes. Decoupling the two closes the ``ran the real
benchmark many times and registered only the best'' attack---the real-rail design is frozen ahead of
the real-rail data, exactly as this method was frozen ahead of any real-rail data.

\paragraph{Verification.} The released manifest verifies every input file by sha256; the harness
re-derives the mock-pipeline verification hash byte-for-byte from the frozen fixtures (no wall-clock is
written into the report). \textbf{This verification hash certifies pipeline reproduction; it is not a
real-rail ranking result.} Only version~1 of this preprint is the canonical frozen artefact; later
revisions, if any, are post-freeze errata.

\section{Related work}
Adjacent efforts touch this ground---registries of payment-capable services, multi-rail routing and
analytics layers, and service-discovery surfaces. PayBench's distinct contribution is
\emph{methodological}: a pre-registered, comparative, settlement-finality benchmark published from a
non-custodial, rail-neutral posture, with the design fixed cryptographically before any scored run. A
neutral, sourced, dated related-work treatment naming specific efforts belongs in the camera-ready
version; it is omitted from this preprint to keep the instrument neutral.

\section{Conclusion}
PayBench operationalises the dominant rail-DX question for agent-to-agent commerce---when can the agent
rely on a payment?---as a comparative, pre-registered benchmark, with a finality doctrine that is honest
about cross-rail trust asymmetry and a statistical method that is honest about its own degeneracies.
The methodology and a reproducible calibrated-mock harness are released now; real-rail rankings follow,
under a separate data pre-registration, when the publication posture permits.

\paragraph{Availability.} Methodology of record, harness, calibration provenance, and the
pre-registration manifest are in the project repository at the frozen commit referenced by the signed
tag \texttt{paybench-prereg-v1.1}. The manifest hash committed by the anchors is
\texttt{sha256:f0b9b079\ldots009d99} (full value in the repository manifest).

\end{document}
